\documentclass[11pt]{article}

\usepackage[margin=1in]{geometry}
\usepackage{amsmath,amssymb,bm}
\usepackage{booktabs}

\title{Small-Strain Flexoelectric Pure Bending with Generalized Traction}
\author{}
\date{}

\newcommand{\dd}{\mathrm{d}}
\newcommand{\eps}{\varepsilon}

\begin{document}
\maketitle

\section{Setting}

Consider a beam in small-strain pure bending.  The reference thickness coordinate is
\[
  Y \in [-a,a], \qquad a=\frac{H}{2}.
\]
The imposed axial bending strain is taken as
\[
  \eps_{11}(Y)=-\kappa Y,
\]
and the transverse strain is left unknown:
\[
  \eps_{22}(Y)=v(Y).
\]
This is the important difference from the simpler plane-stress-style derivation: the
local transverse stress is not prescribed to be zero pointwise.  Instead, the free
surface condition is imposed on the generalized traction.

For an isotropic linear elastic solid,
\[
  C_L=\lambda+2\mu, \qquad C_T=\lambda,
\]
so the local transverse stress is
\begin{equation}
  \sigma_{22}=C_L v-C_T\kappa Y .
  \label{eq:sigma22}
\end{equation}

\section{Open-Circuit Electric Field}

Let a prime denote differentiation with respect to \(Y\).  The one-dimensional
open-circuit electric displacement relation is written as
\begin{equation}
  D_2=\epsilon E_2+\mu_L v'-\mu_T\kappa .
  \label{eq:D2}
\end{equation}
For the open-circuit through-thickness response,
\[
  D_2=0,
\]
and therefore
\begin{equation}
  E_2(Y)=\frac{\mu_T\kappa-\mu_L v'(Y)}{\epsilon}.
  \label{eq:E2-eliminate}
\end{equation}

\section{Higher-Order Stress After Elimination}

Before eliminating the electric field, the relevant higher-order stress component is
\begin{equation}
  S_{222}=H_L v'-H_T\kappa-\mu_L E_2 .
\end{equation}
Using Eq.~\eqref{eq:E2-eliminate},
\begin{equation}
  S_{222}=G_L v'-G_T\kappa ,
  \label{eq:S222}
\end{equation}
where
\begin{equation}
  G_L=H_L+\frac{\mu_L^2}{\epsilon},
  \qquad
  G_T=H_T+\frac{\mu_L\mu_T}{\epsilon}.
  \label{eq:G-general}
\end{equation}

The coefficients \(H_L,H_T\) depend on the strain-gradient/flexoelectric material
form used in the numerical model.  Two useful cases are:
\[
\begin{array}{lll}
\toprule
\text{model} & H_L,\ H_T & G_L,\ G_T \\
\midrule
\text{bare} &
C_L\ell^2,\ C_T\ell^2 &
C_L\ell^2+\mu_L^2/\epsilon,\ C_T\ell^2+\mu_L\mu_T/\epsilon
\\[2mm]
\text{transformed} &
C_L\ell^2-\dfrac{\mu_L^2}{\epsilon-\epsilon_0},\
C_T\ell^2-\dfrac{\mu_L\mu_T}{\epsilon-\epsilon_0}
&
C_L\ell^2-\dfrac{\epsilon_0\mu_L^2}{\epsilon(\epsilon-\epsilon_0)},\
C_T\ell^2-\dfrac{\epsilon_0\mu_L\mu_T}{\epsilon(\epsilon-\epsilon_0)}
\\
\bottomrule
\end{array}
\]
The plotting script uses the bare model by default.  This matches a numerical run
where the additional \(\bar H\) flexoelectric correction is not subtracted.  The
transformed model is kept as an option for runs where that correction is included.

\section{Generalized Free-Surface Conditions}

The through-thickness equilibrium and ordinary traction condition are
\[
  \left(\sigma_{22}-S_{222}'\right)'=0,
  \qquad
  \left.\left(\sigma_{22}-S_{222}'\right)\right|_{Y=\pm a}=0 .
\]
Thus,
\begin{equation}
  \sigma_{22}-S_{222}'=0 .
  \label{eq:generalized-traction}
\end{equation}
Using Eqs.~\eqref{eq:sigma22} and \eqref{eq:S222}, this becomes
\begin{equation}
  C_L v-C_T\kappa Y-G_L v''=0 .
\end{equation}
Equivalently,
\begin{equation}
  v-\xi^2 v''=\beta\kappa Y,
  \qquad
  \xi=\sqrt{\frac{G_L}{C_L}},
  \qquad
  \beta=\frac{C_T}{C_L}.
  \label{eq:ode}
\end{equation}

The higher-order traction is also free at the top and bottom surfaces:
\[
  S_{222}(\pm a)=0.
\]
Therefore,
\begin{equation}
  v'(\pm a)=\gamma\kappa,
  \qquad
  \gamma=\frac{G_T}{G_L}.
  \label{eq:higher-bc}
\end{equation}
Equations~\eqref{eq:ode} and \eqref{eq:higher-bc} are the replacement for the
pointwise assumption \(\sigma_{22}=0\).

\section{Closed Form}

Define
\[
  \eta=\gamma-\beta=\frac{G_T}{G_L}-\frac{C_T}{C_L}.
\]
The solution of Eq.~\eqref{eq:ode} satisfying Eq.~\eqref{eq:higher-bc} is
\begin{equation}
  v(Y)=\beta\kappa Y
  +\kappa\xi\eta\,
  \frac{\sinh(Y/\xi)}{\cosh(a/\xi)} ,
  \label{eq:v}
\end{equation}
and
\begin{equation}
  v'(Y)=\beta\kappa
  +\kappa\eta\,
  \frac{\cosh(Y/\xi)}{\cosh(a/\xi)} .
  \label{eq:vprime}
\end{equation}

The local transverse stress is then
\[
  \sigma_{22}
  =C_L\kappa\xi\eta\,
  \frac{\sinh(Y/\xi)}{\cosh(a/\xi)} .
\]
It is generally not zero.  However, this is exactly balanced by \(S_{222}'\), so
the generalized traction \(\sigma_{22}-S_{222}'\) vanishes.

\section{Electric Potential}

Substitution of Eq.~\eqref{eq:vprime} into Eq.~\eqref{eq:E2-eliminate} gives
\begin{equation}
  E_2(Y)=
  \frac{\kappa}{\epsilon}
  \left[
    \mu_T-\mu_L\beta
    -\mu_L\eta\,
    \frac{\cosh(Y/\xi)}{\cosh(a/\xi)}
  \right].
  \label{eq:E2}
\end{equation}
Let
\[
  A_0=\frac{\kappa}{\epsilon}\left(\mu_T-\mu_L\beta\right),
  \qquad
  B_0=\frac{\mu_L\kappa\xi\eta}{\epsilon}.
\]
Since \(E_2=-\phi'\), and choosing the bottom surface gauge
\[
  \phi(-a)=0,
\]
the electric potential is
\begin{equation}
  \boxed{
  \phi(Y)=
  -A_0(Y+a)
  +B_0
  \left[
    \frac{\sinh(Y/\xi)}{\cosh(a/\xi)}
    +\tanh(a/\xi)
  \right]
  }.
  \label{eq:phi}
\end{equation}

When \(G_L=0\), the boundary-layer length \(\xi\) disappears and the reduced local
solution is
\[
  v(Y)=\beta\kappa Y,\qquad
  E_2=A_0,\qquad
  \phi(Y)=-A_0(Y+a).
\]
If \(G_L<0\), the boundary-layer length is not real, so this closed form is not an
admissible real-valued solution.

\section{Implementation Notes}

The plotting script implements Eq.~\eqref{eq:phi}.  It reads
\(H,L,E,\nu,\ell,\epsilon,\mu_L,\mu_T\) from \texttt{log.txt}, extracts
\(\kappa\) from the selected ParaView step, and uses the bottom-surface gauge
\(\phi(-H/2)=0\).  Its default analytical model is \texttt{bare};
\texttt{transformed} and \texttt{local} are available through
\texttt{--analytical-model}.

\end{document}
